\documentclass[12pt,a4paper,twoside,openright]{book}

% Package imports
\usepackage[utf8]{inputenc}
\usepackage[italian]{babel}
\usepackage[T1]{fontenc}
\usepackage{geometry}
\usepackage{fancyhdr}
\usepackage{titlesec}
\usepackage{tocloft}
\usepackage{setspace}

% Page setup
\geometry{
	top=3cm,
	bottom=3cm,
	left=3.5cm,
	right=2.5cm,
	bindingoffset=0.5cm
}

% Header configuration
\setlength{\headheight}{15pt}
\pagestyle{fancy}
\fancyhf{}
\fancyhead[LE,RO]{\thepage}
\fancyhead[LO]{\rightmark}
\fancyhead[RE]{\leftmark}

% Title formatting
\titleformat{\chapter}[display]
{\normalfont\huge\bfseries}{\chaptertitlename\ \thechapter}{20pt}{\Huge}
\titlespacing*{\chapter}{0pt}{50pt}{40pt}

% TOC formatting
\renewcommand{\cftchapfont}{\bfseries}
\renewcommand{\cftsecfont}{\normalfont}
\renewcommand{\cftsubsecfont}{\normalfont}

% Line spacing
\onehalfspacing

\begin{document}
	
	% Title page
	\begin{titlepage}
		\centering
		\vspace*{1cm}
		
		{\large UNIVERSITÀ STATALE DI MILANO}\\[0.5cm]
		{\large Dipartimento di Informatica "Giovanni Degli Antoni"}\\[0.5cm]
		{\large Corso di Laurea Triennale in Informatica}\\[3cm]
		
		{\Huge\bfseries Sistema di Controllo Avanzato per Robot Dobot CR5:\\[0.5cm]
			Architettura Modulare per Scansione Automatizzata\\[0.5cm]
			con ROS2 e MoveIt2}\\[4cm]
		
		\begin{minipage}{0.4\textwidth}
			\begin{flushleft}
				\large
				\textbf{Relatore:}\\
				Prof.ssa [Nome Relatrice]\\[1cm]
				\textbf{Correlatore:}\\
				Ing. [Nome Correlatore]
			\end{flushleft}
		\end{minipage}
		\hfill
		\begin{minipage}{0.4\textwidth}
			\begin{flushright}
				\large
				\textbf{Candidato:}\\
				Andrea [Cognome]\\
				Matricola: [Numero]
			\end{flushright}
		\end{minipage}
		
		\vfill
		
		{\large Anno Accademico 2024-2025}
		
	\end{titlepage}
	
	% Table of Contents
	\tableofcontents
	\cleardoublepage
	
	% Main content structure with detailed index
	
	\chapter{Introduzione}
	\section{Definizione del Problema}
	\subsection{Limitazioni dei sistemi software standard per robot Dobot}
	\subsection{Necessità di implementazione di un robot autonomo operante in sicurezza}
	
	\section{Obiettivi della Tesi}
	\subsection{Sviluppo dell'infrastruttura base per il controllo del Dobot CR5}
	\subsection{Implementazione di un sistema di comunicazione tra ZED e Dobot}
	\subsection{Progettazione e implementazione di algoritmi per la scansione automatizzata delle piante}
	
	\section{Struttura della Tesi}
	
	\chapter{Stato dell'Arte e Tecnologie}
	\section{Robotics Operating System 2 (ROS2)}
	\subsection{Architettura distribuita e Data Distribution Service (DDS)}
	\subsection{Evoluzione da ROS a ROS2: sicurezza e real-time}
	\subsection{Paradigmi di comunicazione: topics, services e actions}
	\subsection{Gestione del ciclo di vita dei nodi (Lifecycle Management)}
	
	\section{MoveIt2: Framework di Motion Planning}
	\subsection{Architettura del sistema di pianificazione}
	\subsection{Algoritmi di path planning: RRTConnect, EST, PRM}
	\subsection{Collision detection e scene management}
	\subsection{Integrazione con robot hardware tramite control interfaces}
	
	\section{Robot Dobot CR5: Specifiche e Limitazioni}
	\subsection{Caratteristiche hardware: 6 DOF, precisione e workspace}
	\subsection{Protocollo TCP/IP nativo e interfacce disponibili}
	\subsection{Analisi del repository ufficiale DOBOT\_6Axis\_ROS2\_V3}
	\subsection{Identificazione delle lacune funzionali}
	
	\section{Metodologie di Sviluppo Software per Robotica}
	\subsection{Design patterns per sistemi robotici distribuiti}
	\subsection{Architetture modulari e separation of concerns}
	\subsection{Testing e validazione in ambienti simulati}
	
	\chapter{Progettazione del Sistema di Pianificazione e Movimento}
	\section{Analisi del Sistema Base Dobot}
	\subsection{Reverse engineering del repository GitHub ufficiale}
	\subsection{Mappatura delle funzionalità mancanti}
	\subsection{Descrizione delle modifiche necessarie al codice esistente: aggiunta del settimo joint per la rappresentazione della camera e modifica del canale di pubblicazione delle trasformate dei joint per l'integrazione con MoveIt}
	
	\section{Requisiti Funzionali del Sistema Esteso}
	\subsection{Raggiungimento di qualsiasi punto fisicamente accessibile nel workspace}
	\subsection{Gestione dinamica di scene complesse con ostacoli}
	\subsection{Implementazione di un sistema di avvio automatizzato}
	
	\section{Requisiti Non Funzionali (METTERE?)}
	
	
	\chapter{Implementazione dell'Architettura Modulare}
	\section{Architettura del Sistema Proposto}
	\subsection{Modulo di attivazione dei driver del Dobot}
	\subsection{Modulo di lancio dei componenti necessari per l'interazione con MoveIt}
	\subsection{Progettazione modulare con schema architetturale (grafico)}
	
	\section{Package \texttt{dobot\_bringup\_v3}, File \texttt{dobot\_bringup\_ros2.launch.py}: Sistema Core}
	\subsection{Descrizione del funzionamento e integrazione con l'infrastruttura ROS2}
	
	\section{Package \texttt{dobot\_bringup\_v3}, File \texttt{translater.py}}
	\subsection{Analisi del funzionamento e ruolo nell'architettura di controllo}
	
	
	
	\chapter{Progettazione del Sistema di Scansione delle Piante}
	\section{Analisi dei Dati Estraibili dalla ZED}
	\subsection{Tipologie di dati acquisibili dal sensore}
	\subsection{Metodologie per la trasmissione dei dati}
	
	\section{Requisiti Funzionali del Sistema}
	\subsection{Comunicazione bidirezionale con la ZED}
	\subsection{Capacità di elaborazione simultanea di multiple piante}
	\subsection{Ottimizzazione degli algoritmi di scansione}
	
	\section{Requisiti Non Funzionali}
	\subsection{Modularità e riusabilità del codice}
	
	\chapter{Implementazione della Struttura di Scansione}
	\section{Architettura del Sistema Proposto}
	\subsection{Canale di comunicazione con la ZED}
	\subsection{Elaborazione della mappa ricevuta dalla ZED}
	\subsection{Processo di scansione utilizzando la mappa rielaborata}
	\subsection{Progettazione modulare con schema architetturale (grafico)}
	
	\section{Package \texttt{dobot\_bringup\_v3}, File \texttt{zed\_listener.py}}
	\subsection{Analisi del funzionamento}
	\subsection{Struttura dei messaggi condivisi: acquisizione dalla ZED e ricezione da parte del Dobot}
	
	\section{Package \texttt{dobot\_bringup\_v3}, File \texttt{reworked\_map\_node.py}}
	\subsection{Sottoscrizione al topic Boing per la ricezione della mappa inviata dalla ZED}
	\subsection{Generazione di circonferenze multiple basate sull'altezza della pianta}
	\subsection{Elaborazione dei punti di ogni circonferenza per l'esclusione di punti irraggiungibili}
	\subsection{Modifica del messaggio ZED e pubblicazione sul topic reworked\_map}
	
	\section{Package \texttt{cr5\_moveit\_cpp\_demo}, File \texttt{move\_cr5\_node\_modular.cpp}}
	\subsection{Responsabilità della classe: coordinamento con PlantProcessor per l'avvio della scansione}
	\subsection{Sottoscrizione al topic reworked\_map per la ricezione della mappa rielaborata}
	\subsection{Inizializzazione dei parametri di sistema}
	
	\section{Package \texttt{cr5\_moveit\_cpp\_demo}, File \texttt{plant\_processor.cpp}}
	\subsection{Responsabilità della classe e interfacciamento con i moduli SceneManager e PlantScanner}
	\subsection{Ricezione della mappa e invocazione del SceneManager per l'aggiunta di ostacoli e piante}
	\subsection{Elaborazione sequenziale di ogni pianta}
	\subsection{Elaborazione sequenziale di ogni circonferenza tramite invocazione del PlantScanner}
	
	\section{Package \texttt{cr5\_moveit\_cpp\_demo}, File \texttt{scene\_manager.cpp}}
	\subsection{Responsabilità della classe e interfacciamento con Planning Interface di MoveIt per l'aggiunta di ostacoli}
	\subsection{Gestione dell'ambiente: aggiunta di pavimento, piante e ostacoli}
	
	\section{Package \texttt{cr5\_moveit\_cpp\_demo}, File \texttt{plant\_scanner.cpp}}
	\subsection{Responsabilità della classe e interfacciamento con il modulo MovementController}
	\subsection{Identificazione degli algoritmi per i punti di confine}
	\subsection{Tentativo di raggiungimento dei punti di confine tramite invocazione del MovementController}
	\subsection{Gestione della copertura: esclusione dei punti vicini basata sul campo visivo della camera}
	
	\section{Package \texttt{cr5\_moveit\_cpp\_demo}, File \texttt{movement\_controller.cpp}}
	\subsection{Responsabilità della classe e interfacciamento con Planning Interface e MoveGroup di MoveIt per l'attuazione del movimento}
	\subsection{Gestione dei casi di successo e fallimento nell'esecuzione dei movimenti}
	
	
	\chapter{Sistema di Gestione e Lancio dei Moduli}
	
	\section{Configurazione per Dobot Reale}
	\subsection{Package \texttt{cr5\_moveit\_cpp\_demo}, File \texttt{./dobot\_startup.sh}}
	\subsubsection{Funzionalità e moduli di sistema avviati}
	
	\subsection{Package \texttt{cr5\_moveit}, File \texttt{full\_bringup.launch.py}}
	\subsubsection{Architettura di avvio e componenti lanciati}
	
	\subsection{Package \texttt{cr5\_moveit\_cpp\_demo}, File \texttt{./dobot\_stop.sh}}
	\subsubsection{Procedura di arresto sicuro del sistema}
	
	\section{Configurazione per Dobot Simulato}
	\subsection{Package \texttt{dobot\_moveit}, File \texttt{dobot\_moveit.launch.py}}
	\subsubsection{Architettura di simulazione e componenti del sistema virtuale}
	
	\subsection{Package \texttt{cr5\_moveit\_cpp\_demo}, File \texttt{scripts\_launch.launch.py}}
	\subsubsection{Coordinamento dei moduli in ambiente simulato}
	
	\subsection{Package \texttt{cr5\_moveit\_cpp\_demo}, File \texttt{fake\_zed.cpp}}
	\subsubsection{Implementazione del simulatore ZED per test di sistema}
	
	\subsection{Package \texttt{cr5\_moveit\_cpp\_demo}, File \texttt{move\_cr5\_node\_modular.cpp}}
	\subsubsection{Modalità di esecuzione in ambiente reale tramite script di avvio}
	
	
	\chapter{Attività di Testing e Validazione del Sistema}
	\section{Testing in Ambiente Simulato}
	\subsection{Obiettivi della validazione in simulazione}
	\subsubsection{Test con configurazioni multiple di piante}
	\subsection{Risultati e osservazioni critiche}
	\subsubsection{Analisi delle difficoltà nella scansione completa delle piante}
	
	\section{Testing con Sistema Dobot Reale}
	\subsection{Obiettivi della validazione in ambiente reale}
	\subsubsection{Verifica della corrispondenza tra simulazione e realtà operativa}
	\subsection{Risultati e confronto prestazionale}
	\subsubsection{Comportamento sostanzialmente identico con eccezioni marginali}
	\subsubsection{Analisi dei casi di fallimento nella raggiungibilità dei punti}
	
	
	\chapter{Conclusioni e Sviluppi Futuri}
	\section{Sintesi dei Risultati}
	\subsection{Valutazione complessiva dei risultati ottenuti}
	\subsubsection{Soddisfacimento degli obiettivi prefissati}
	\subsubsection{Identificazione di problematiche risolvibili mediante applicazioni alternative}
	
	\section{Prospettive di Sviluppo Futuro}
	\subsection{Integrazione con sistemi robotici mobili}
	\subsubsection{Implementazione di capacità di spostamento autonomo del sistema}
	\subsection{Progettazione orientata all'adattabilità}
	\subsubsection{Architettura modulare per facile adattamento a nuove configurazioni spaziali}
	\subsection{Estensibilità applicativa}
	\subsubsection{Adattabilità del sistema a diverse tipologie di applicazioni robotiche industriali}
	

	% Bibliografia
	\chapter*{Bibliografia}
	\addcontentsline{toc}{chapter}{Bibliografia}
	
	% Indice delle figure
	\listoffigures
	
	% Indice delle tabelle
	\listoftables
	
\end{document}